\scanpage{7}
\begin{center}- 179 -\end{center}
\margin{n° 391}

conséquences sur les $m_{ij}$ qui permettront de dresser la liste de toutes les possibilités pour les schémas de Coxeter correspondants ; il restera ensuite à voir que ces possibilités sont effectivement réalisées, c'est-à-dire que les groupes correspondants $G(\underline M)$ sont \uline{finis}, (ou encore que la forme $q$ est positive non dégénérée).

LEMME 1.- \uline{Pour tout} $i$, \uline{on a} $\displaystyle\sum_{j\ne i}q_{ij}^2<1$.

On peut se borner à ne considérer que les indices $j$ pour lesquels $q_{ij}\ne0$, c'est-à-dire ceux qui sont liés à $i$ dans $X$. Si $j$, $j'$ sont deux de ces indices distincts, $\{j,j'\}$ n'est pas une arête, sinon $\{i,j\}$, $\{j,j'\}$ et $\{j',i\}$ formeraient un circuit dans $X$ ; donc $e_j$ et $e_{j'}$ sont \uline{orthogonaux}, et le lemme résulte du fait que, dans $E$, la distance $d$ de $e_i$ au sous-espace vectoriel $F$ engendré par les $e_j$, ($j$ lié à $i$), est donnée par $\displaystyle d^2=1-\sum_j(e_i\mid e_j)^2=1-\sum_{j\ne i}q_{ij}^2$ en vertu du th. de Pythagore, les $e_j$ formant une base orthonormale de $F$.

COROLLAIRE 1.- \uline{Un sommet de} $X$ \uline{ne peut appartenir qu'à} 3 \uline{arêtes au plus}.

En effet, si $i$ est lié à $h$ autres sommets, les relations $q_{ij}^2\geq1/4$ pour ces sommets $j$ entraînent $h/4<1$ par le lemme 1, donc $h\leq3$.

COROLLAIRE 2.- \uline{Si} $i$ \uline{appartient à} 3 \uline{arêtes}, \uline{ces arêtes sont d'ordre} 3.

Comme $\cos(\pi/m)\geq\sqrt2/2$ pour $m\geq4$, (\S\ 1, n° 3, prop. 8), on aurait en effet, dans le cas contraire, $\displaystyle\sum_{j\ne i}q_{ij}^2\geq\frac14+\frac14+\frac12=1$, ce qui contredit le lemme 1.

COROLLAIRE 3.- \uline{S'il existe une arête d'ordre} $\geq6$, \uline{on a} $\ell=2$.

Soit en effet $\{i,j\}$ cette arête. Si on avait $\ell>2$, un des sommets $i$, $j$, (par exemple $i$), serait lié à un troisième sommet $j'$, puisque $X$ est connexe.
\mspage{6}
